\section{Revision and source of the report}
This revision responds to the complete R30 report on manuscript commit \texttt{3394815}. The report is preserved on review commit \texttt{e9bc144}. We read that report together with the intervening R31/R32 construction and numerical records, rather than treating the reviewed R30 manuscript as the latest available science. R34 is based on frozen parent \texttt{88015b7}, containing the executed R32 package and staged R33 sources and is confined to a new revision branch. The R32 source and executed evidence are inherited work, explicitly identified below; the unrestricted-cost study is freshly executed in R34, and its 252 base certificates are supplemented by 42 new restart-propagated certificates. The full 294-object audit separates inherited evidence from newly executed results.

The substantive addition is a constructive global intervention-cost interval together with a backward propagation that exploits simultaneous feasibility at every restart. Theorem~\ref{art-thm:restart} proves the propagated bound for randomized Markov policies and the precisely specified conditional-history extension; it is not only an inner-class optimality statement. A computed lower operating witness, unrestricted scalarized Bellman values, and an independently feasible deployment bound the true minimum implementation cost beyond the conservative local action class. The revised manuscript retains every previous theorem, its original stopped-economy target, the complete boundary proof, all adverse comparisons, and the earlier sensitivity frontier. The historical annex includes the entire R32 article, supplement, response, computation record, and its existing historical annex. No earlier negative experiment has been discarded.

\section{Main scientific findings}
\subsection{F1: Constructing and charging the witnesses}
The constructive upper-witness algorithm and endogenous-cost computation introduced in R32 are retained in Sections~\ref{art-sec:construction} and \ref{art-sec:revision}. R34 adds the computed lower witness $L=U-b$ from actual checked compression defects, rather than taking an optimal value as input. Section~\ref{art-sec:global} constructs all support values, their policies, and their operating and implementation values. Their costs are included. The exact operating dynamic program is run separately for comparison, not used as a supplied certificate input.

The staged R33 sources identify an earlier timing omission. R34 uses their complete-cost constructor from the first new execution: the class candidate is rebuilt, checked for exact equality to the inherited object, and charged for witness construction, action masks, cost recursion, evaluation, encoding, and checking. The source archive remains unchanged; we do not claim a new execution of an unavailable prior attempt. Table~\ref{art-tab:global-work} and the complete timing JSON distinguish individual costs, shared work, and the additional final audit.

\subsection{F2: Nontrivial Bellman coupling}
The R32 continuous-state maintenance economy remains the primary computational application. Maintain and replace have lower immediate reward than defer. They can be justified only by their action-dependent future condition distributions. The unrestricted R34 support problems retain those same transitions, operating returns, and state-dependent revision exposure. No continuation cancellation or primitive reward-sign policy is introduced. The auxiliary action-independent R30 examples are retained as historical controls, not as the principal coupled computation.

\subsection{F3: Strong controls and the value of additional work}
We retain exact structural dynamic programming and both spline controls with the same information. The R34 class comparator is reconstructed from primitives. All 18 installed-neural configurations obtain strictly lower implementation cost outside that conservative class while satisfying their unchanged operating tolerances and preserving their installed operating values. This is an economic improvement over the restricted-cost solution, not evidence that the more expensive portfolio is faster than exact operating dynamic programming. The portfolio delivers an additional output: a lower cost bound against every globally feasible policy.

For the most informative $T=12$, $.05$, seed-31003 comparison, cost falls from approximately $1.404696$ to $.674484$, but the unrestricted lower bound is approximately $.041222$. The reported $.633262$ global gap makes the remaining economic uncertainty explicit. We do not replace the referee's computational-advantage question with a false speedup claim. A verified end-to-end advantage over the best classical solver in a difficult higher-dimensional application is not established by this cohort.

\subsection{F4: Endogenous occupancy}
The inherited cost Bellman recursion evaluates future interventions under the revised policy's own transitions. Its six predeclared stress configurations have positive dynamic-versus-pointwise savings, while the other 36 comparisons have zero savings; all are retained. R34 goes further by optimizing unrestricted support policies under their own transition laws, then independently recomputing their complete implementation costs. Neither policy selection nor the class-free lower bound substitutes the installed occupancy measure for the deployed one.

\subsection{F5: State cardinality, dimension, and geometry}
No trillion-state headline is used. The current contribution is a finite whole-interval representation for a one-dimensional action-dependent economy. Section~\ref{art-sec:work} records the generated representations and integer bit lengths; Table~\ref{art-tab:global-work} adds the support-value geometry and peak process memory. The output-sensitive work bound is not a dimension-independent rate. Earlier $2^{40}$ statements survive only in their historical context as representation observations, not as a run of a trillion-state dynamic program.

\subsection{F6: The role of neural proposals}
We preserve neural proposals as installed policies whose revision has an economic cost, and keep the exact distinction between trained coefficients, rational compiled policies, and hardware execution. The experiment has three inherited neural seeds, not a newly sampled large cohort. All 18 such cases improve relative to the conservative class, while all twelve spline cases have zero-cost globally feasible incumbents. These findings support an installed-policy revision method that accommodates neural rules, not neural-specific acceleration or superiority over the spline controls. No sampled fitting loss is promoted to a verified whole-state approximation bound.

\subsection{F7: Complexity and a quantitative proposal condition}
The constructive compression bound, charged output-sensitive complexity, and verified proposal-margin condition in Section~\ref{art-sec:work} are retained. Section~\ref{art-sec:global} adds the work of every multiplier solve, its own-policy evaluation, lower-envelope construction, encoding, and checking. The largest declared multiplier has a model-based feasibility guarantee; feasibility is not selected after inspecting successful cases. Certificate size is a measured construction output and can grow with horizon, geometry, and bit length. No favorable unmeasured high-dimensional scaling law is assumed.

\subsection{F8: The unchanged original all-domain target}
The original stopped consumption--portfolio primitives and the exact target remain in Section~\ref{art-sec:stopped}. The retained all-domain bound $7.181834580823298$ is not replaced by the maintenance-model bounds, and the missing 47-coordinate current-state bridge is not declared complete. The interval-wide constant-control result is a concrete advance on an explicitly identified restriction of that original economy. R34's new cost bounds are for the coupled maintenance application, not a disguised certificate for the full stopped feedback problem. The research objective and technical material are preserved without claiming a numerical reduction that has not been computed.

\subsection{F9: Completing the active-boundary scalar problem}
The full R32 global scalar optimization and its proof are retained, rather than returning to the five-point diagnostic criticized in R30. The proof establishes a positive derivative throughout both halves of the constant-adjustment interval, bounds the distant-boundary contribution, and identifies the unique maximizer $\theta=1/5$ for the declared active-boundary restart and fixed consumption/portfolio controls. The supplement retains all rational details. This is a complete scalar-class result, not an unrestricted feedback-control certificate. Its proof and endpoint enclosure are identified as inherited work; R34 reruns the elementary exact-rational inequality checks, but does not claim a new endpoint quadrature or a new feedback certificate.

\subsection{F10: Economic units and sensitivity}
No universal 60/60 success count at an arbitrary external revision price is reinstated. We retain the entire named-policy sensitivity frontier and normalized intervention units. R34 reports complete cost lower/upper intervals at the two fixed operating tolerances. Its eight multipliers are computational scalarization parameters, not calibrated monetary estimates and not the earlier illustrative external welfare price. The new economic conclusions are cost reductions and certified global gaps in normalized implementation units, not a calibration claim.

\subsection{F11: The theorem-level contribution}
The manuscript explicitly credits classical Bellman comparison and weak duality. The new theorem is not presented as a new duality principle. It composes constructively computed two-sided operating witnesses with unrestricted support certificates, a globally feasible deployment, and a checked global implementation-cost interval for the simultaneous all-state, all-restart operating constraint. This differs from asserting optimality only after a conservative eligible-action set is supplied. It also distinguishes finite policy selection from a lower bound valid for all feasible policies, and does not assume strong duality or a zero gap.

Theorem~\ref{art-thm:restart} adds a specific mechanism absent from the separate-restart support inequality: feasible future intervention costs are propagated through the actual action-dependent transition law, while a nonnegative local multiplier enforces a necessary current operating inequality. No action is discarded, so feasible randomized mixtures remain covered. The improvement is verified on the same upper deployment, with \RestartTightened{} strict lower-bound gains and all unchanged outcomes retained. Exact min/max proof selectors allow a checker to establish these identities without invoking optimization. The additional output-sensitive work is charged in equation~\eqref{art-eq:restart-work}; no new unrestricted support solve is required.

\subsection{F12: A dominant theorem-to-computation chain}
The main chain is now: coupled maintenance primitives; constructive operating witnesses; local-class repair and its limitations; unrestricted support values; a feasible cost-improving deployment; and a lower bound against the original globally feasible set. The 294 independently checked objects, including the additional restart-propagated lower bounds, instantiate this chain in a single unchanged economic model. The stopped-boundary scalar theorem remains a separate, fully proved advancement with an explicit scope. Its results are not counted as evidence for the maintenance algorithm or used to fill a missing high-dimensional scalability result.

The added global-certificate chain now reaches the all-restart restriction itself rather than only supplying a separate-restart relaxation. For the largest absolute gain, the cost interval tightens from $[\RestartOldLower,\RestartUpper]$ to $[\RestartNewLower,\RestartUpper]$ without changing the operating policy. Positive residual gaps remain visible; this is a strengthened verified numerical statement, not a claim that finite scalarization automatically solves every constrained optimum.

\section{Technical comments and disposition}
\subsection{T1 and T2: Certificate-input work.}
Upper and lower witnesses, all own-policy values, all support solves, and the newly reconstructed class comparator are constructed and charged. The final checker reads serialized objects and is timed separately. The frozen historical execution and staged source archive remain available; the complete-cost R34 executions and the additional propagation costs have their own ledgers.

\subsection{T3: Deployment amortization.}
The inherited break-even query formula and its assumptions remain in Section~\ref{art-sec:work}. No measured break-even is claimed because online-query and loading costs have not been benchmarked under an agreed deployment implementation. The new offline computation is not represented as an online amortization experiment. This requested deployment benchmark remains unmeasured, rather than being claimed complete from a formula alone.

\subsection{T4: Dimension.}
The domain dimension is one throughout the new computation. We report certificate size and process memory without substituting state-cardinality independence for dimension independence.

\subsection{T5: Occupancy.}
Both the inherited class-cost recursion and the new unrestricted support recursion use the revised action's transition kernel. Complete exact own-policy evaluations check the resulting implementation costs.

\subsection{T6: Proposal quality.}
The verified $\rho+2\kappa$ deficit bound and its zero-revision implication are preserved. The new global feasibility bound depends on bounded implementation cost and computed witness defects, not on an unverified neural approximation estimate. The observed neural-case improvements do not establish a favorable approximation-rate comparison against splines.

\subsection{T7: Deployment semantics.}
Certificates cover rational compiled policies, including isolated tie points. Floating-point hardware neural execution is not silently included. Changed endpoint and action values are explicitly exercised in adversarial checker tests.

\subsection{T8: Units.}
One cost unit is a normalized implementation intervention weighted by $1+x$. Multipliers used in the support computation are distinguished from externally calibrated economic prices. No empirical dollar interpretation is introduced.

\subsection{T9: Boundary optimization.}
The complete interval-wide scalar proof and its rational bounds are retained in the article and supplement. We distinguish its unique scalar optimum from the still separate all-domain feedback objective.

\subsection{T10: Preservation and editorial focus.}
All previous technical content remains in the revision or the lossless historical annex. New global-cost theorems and the coupled application are the dominant additions. The preservation map records the unchanged source hashes, deterministic text transformations, and inserted sections. Earlier responses and computations remain accessible without forcing the reader to infer the new contribution from a sequence of response letters.

\section{Review package and reproducibility}
The root article, supplement, response, computational record, and historical annex are \texttt{ECTA\_R34}, \texttt{SUPP\_R34}, \texttt{RESPONSE\_R34}, \texttt{COMPUTATION\_R34}, and \texttt{HISTORY\_R34}. Each has an English LaTeX source; the principal documents use the repository's unmodified official Econometric Society class. The new revision directory contains the fixed protocol, implementation and provenance notes, readable constructors and independent checkers, all 42 outcome rows and 168 multiplier rows, all 294 exact certificates, timing records, mutation tests, source hashes, and publication manifest. The existing review and revision branches are not modified.

\section{Additional review-ready check}
The complete 42-case propagation table and the 12 corruption tests are in the supplement and execution records. Four-state finite-tree checks include randomized mixtures and explicitly record the eligible count, rather than checking only actionwise feasibility. The current protocol fixes the propagation before its implementation, while identifying the support outcomes that were already known. All newly generated sources, proof objects, tables, and PDFs are committed on the R34 branch; the source report and every preceding branch remain unchanged.
